% RQ4 Conclusion: Treasury Resilience

\subsection{Summary}

We have presented a systematic analysis of treasury management policies through multi-agent simulation. Across \experimentruns{} simulation runs testing combinations of buffer fractions, spending limits, and emergency top-up mechanisms, we characterized the resilience-efficiency tradeoff in DAO treasury management.

Key findings:

\begin{enumerate}
    \item \textbf{Buffers are essential}: 10-20\% buffer fractions substantially reduce treasury volatility without excessive capital lockup. No buffer leaves treasuries vulnerable to shocks.

    \item \textbf{Spending limits extend runway}: 2-5\% periodic limits prevent treasury depletion but create operational constraints. Limits should be calibrated to sustainable burn rate.

    \item \textbf{Emergency mechanisms provide backstop}: Top-up mechanisms reduce maximum drawdown severity. Revenue diversion is gentler; spending freeze is more aggressive.

    \item \textbf{Market regime matters}: Buffer effectiveness varies by market conditions. Higher buffers are most valuable during bear markets when volatility peaks.
\end{enumerate}

\subsection{Contributions}

This paper contributes:

\begin{enumerate}
    \item \textbf{Resilience metrics} for treasury health assessment
    \item \textbf{Parameter guidelines} for buffer fractions, spending limits, and top-up mechanisms
    \item \textbf{Policy design framework} with tiered recommendations
    \item \textbf{Market-adaptive strategies} for different risk profiles
\end{enumerate}

\subsection{Practical Recommendations}

For DAO practitioners designing treasury policies:

\begin{enumerate}
    \item \textbf{Establish explicit buffer}: Target 15-20\% of treasury as reserve
    \item \textbf{Set spending limits}: Cap periodic outflows at sustainable rate (consider 5\% of treasury per quarter)
    \item \textbf{Define top-up triggers}: Activate emergency mechanisms when buffer falls below 50\% of target
    \item \textbf{Monitor and adapt}: Track volatility, drawdown, and runway metrics; adjust policies as conditions change
    \item \textbf{Diversify holdings}: Reduce native token concentration to limit correlated risk
\end{enumerate}

\subsection{Implications}

For DAO treasuries, our results underscore that ad hoc management is insufficient. The billions held in DAO treasuries deserve systematic policy frameworks comparable to corporate or institutional treasury management.

For researchers, our framework enables continued investigation of treasury dynamics. Integration with protocol economics, market microstructure, and strategic behavior offers rich avenues for future work.

\subsection{Closing Remarks}

Treasury resilience is not glamorous but it is foundational. DAOs cannot pursue their missions---development, grants, community building---without fiscal sustainability. The policies we analyze are not exotic innovations but straightforward adaptations of proven financial practices.

We release our simulation framework to enable DAOs to model their specific treasury dynamics and design policies appropriate to their context and risk tolerance.
